
\section{Implementation Details}
The project is implemented in Python using OpenCV, NumPy, scikit-learn, SciPy, Pillow, Pandas, and Matplotlib. The main program is menu-driven and saves every output to a predefined folder, so the results are reproducible and do not remain only as screen previews.

\begin{table}[H]
\centering
\caption{Main program components.}
\begin{tabular}{@{}p{0.28\textwidth}p{0.64\textwidth}@{}}
\toprule
File & Purpose \\
\midrule
\texttt{main\_3\_final.py} & Builds feature spaces, runs MiniBatch K-means and Mean Shift, saves raw and post-processed segmentation outputs. \\
\texttt{datasetCreation.py} & Converts annotation JSON polygons into mask, contour, and contour-overlay GT images. \\
\texttt{evaluate\_2.py} & Evaluates segmentation masks against GT using label matching, confusion matrices, Precision, Recall, IoU, and Boundary F1-score. \\
\texttt{fast\_eval\_minimal.py} & Helper script used in this package to regenerate compact CSV summaries quickly. \\
\texttt{evaluate\_peer\_selected.py} & Helper script used to evaluate selected best predictions against peer GT. \\
\bottomrule
\end{tabular}
\label{tab:program-components}
\end{table}

\subsection{Output folder convention}
The output folder is organized by image, feature space, and algorithm. For example, a typical result is stored as:
\begin{lstlisting}
outputs/Img01/xy_hsv/Half Resize PCA 95 Mean Shift/Post Processed Img01-meanshift-0.8.png
\end{lstlisting}
This path tells us the image ID, feature set, Mean Shift computation mode, and bandwidth parameter. K-means outputs are stored similarly under a \texttt{KMeans} subfolder.

\subsection{Evaluation protocol}
The machine-generated segmentation labels do not necessarily have the same numeric/color IDs as the GT labels, so label matching is required before measuring metrics. This package uses many-to-one label matching in the regenerated CSV summaries. That choice is suitable for segmentation because an algorithm may over-segment a single GT object into several pieces; those pieces should still be allowed to map to the same GT class for pixel-level scoring. The normalized confusion matrix is then computed after matching.

The per-class measures are:
\begin{align}
\mathrm{Precision} &= \frac{TP}{TP + FP}, &
\mathrm{Recall} &= \frac{TP}{TP + FN}, &
\mathrm{IoU} &= \frac{TP}{TP + FP + FN}.
\end{align}
Boundary F1-score is computed by extracting class boundaries and allowing a tolerance of 2 pixels. Macro scores are averaged over foreground classes, with background excluded from the macro average but still recorded in the per-class CSV files.
